\documentclass[UTF8,fontset=fandol,11pt]{ctexart}
\usepackage[a4paper,margin=25mm]{geometry}
\usepackage{amsmath,amssymb,amsthm,mathtools}
\usepackage{xcolor}
\usepackage{booktabs,array,enumitem,fancyhdr,hyperref}
\hypersetup{colorlinks=true,linkcolor=blue!50!black,urlcolor=blue!50!black}
\setlength{\emergencystretch}{2em}
\setlength{\parskip}{0.3em}
\linespread{1.16}
\allowdisplaybreaks
\newtheorem{theorem}{定理}[section]
\newtheorem{proposition}[theorem]{命题}
\newtheorem{lemma}[theorem]{引理}
\theoremstyle{definition}
\newtheorem{definition}[theorem]{定义}
\theoremstyle{remark}
\newtheorem{remark}[theorem]{说明}
\newcommand{\dd}{\,\mathrm d}
\newcommand{\ii}{\mathrm i}
\newcommand{\Rea}{\operatorname{Re}}
\newcommand{\Ima}{\operatorname{Im}}
\newcommand{\Li}{\operatorname{Li}}
\newcommand{\supp}{\operatorname{supp}}
\pagestyle{fancy}
\fancyhf{}
\fancyhead[L]{\small 学习笔记}
\fancyhead[R]{\small 初稿}
\fancyfoot[C]{\thepage}
\setlength{\headheight}{15pt}
\author{}
\date{2026年10月22日初稿}
\makeatletter
\renewcommand{\maketitle}{%
  \thispagestyle{plain}%
  \begin{center}
    {\LARGE\bfseries \@title\par}
    \vspace{0.6em}
    {\small \@date\par}
  \end{center}
  \vspace{0.6em}}
\makeatother
\newcommand{\draftnotice}{\noindent{\small
本文为供逐篇审阅的学习笔记初稿，不是新增研究成果。
文中区分直接证明与引用结果；正文不依赖本文。
本轮已作公式、条件及排版初步核查，尚待作者逐句审定。}\par\medskip}

\title{怎样证明一次数值计算可信}
\begin{document}
\maketitle
\draftnotice
\section{近似值需要回答哪些问题}
计算器给出一个很小的残差，并不自动证明附近存在根，也不自动给出根的位置误差。
应分清四件事：函数值计算是否可靠，所给区域内是否存在根，根是否唯一，
以及找到的根是否覆盖所研究范围中的全部根。
本篇从 $\sqrt2$ 的实根开始，给出可逐项核查的认证，再说明复根计数多出的任务。
重点是可以手工检查的小型证书；区间运算的软件实现及形式化证明的信任基础，
另见学习笔记第 16 篇。

\section{二分法给出的是区间证明}
令 $f(x)=x^2-2$。$f(1)<0<f(2)$，连续性保证至少有一根；
$f'(x)=2x>0$（$1\le x\le2$）保证至多一根。
每次二分保留符号相反的半区间。经过 $N$ 次，区间宽为 $2^{-N}$，
根仍在其中；以中点作近似时，误差至多为 $2^{-N-1}$。
这个结论来自区间不变性，不依赖迭代值看起来是否已经稳定。

例如取有理数
\[
 a=\frac{141421356237}{10^{11}},\qquad
 b=\frac{141421356238}{10^{11}}.
\]
用整数运算可核查
$141421356237^2<2\cdot10^{22}<141421356238^2$。
所以 $\sqrt2\in(a,b)$，中点误差小于 $5\times10^{-12}$。
这里的十进制端点是精确有理数，不是未经控制的浮点输出。

\section{小残差何时能控制位置误差}
若根为 $\alpha$，且在包含 $x$ 与 $\alpha$ 的区间内 $f'$ 同号并有 $|f'|\ge m>0$，
由中值定理
\[
 |x-\alpha|\le |f(x)|/m.
\]
对本例在 $[1,2]$ 可取 $m=2$。
若只知道浮点计算 $\widetilde f(x)$，还须有可靠界
$|f(x)-\widetilde f(x)|\le\eta$，才能写
$|x-\alpha|\le(|\widetilde f(x)|+\eta)/m$。
存在性以及 $x,\alpha$ 所在的区间，都必须另有依据。

反例是 $f(x)=\varepsilon(x-1)$：在 $x=0$ 残差仅为 $\varepsilon$，
但根仍距它 $1$。对 $f(x)=(x-1)^2$，残差与根误差是平方关系，
而且导数在根处为零，不能使用固定正下界的上述估计。
这就是“残差小”与“位置误差小”之间需要条件的原因。

\section{Newton 法很快，但仍须认证}
本例的 Newton 迭代为
\[
 x_{n+1}=\frac12(x_n+2/x_n),\qquad x_n>0.
\]
记 $\alpha=\sqrt2$，利用 $\alpha^2=2$ 精确计算
\[
 x_{n+1}-\alpha=\frac{(x_n-\alpha)^2}{2x_n}\ge0.
\]
所以一次迭代后总在根的右侧；若 $x_n\ge\alpha$，
则 $x_{n+1}\le x_n$，正的单调序列收敛，其极限方程给出极限为 $\alpha$。
误差等式还说明根附近的二次收敛。

在实际浮点计算中，等式会受到舍入误差影响。
可采用有理迭代或向外舍入的区间运算，最后再核验一对异号端点。
本文给出的速度分析不能代替某次具体机器运行的误差认证。

\section{复平面上的一个完整小例子}
没有实数的正负号时，可用 Rouché 定理计数。
在圆盘 $|z-c|<r$ 中取 $c=7/5$、$r=1/50$，写
\[
 z^2-2=2c(z-c)+(c^2-2)+(z-c)^2.
\]
圆周上
\[
 |(c^2-2)+(z-c)^2|
 \le\frac1{25}+\frac1{2500}
 =\frac{101}{2500}
 <\frac7{125}=|2c(z-c)|.
\]
由 Rouché 定理，$z^2-2$ 与线性函数 $2c(z-c)$ 在圆盘内有同样的一个零点，计重数。
由于圆盘关于实轴对称且多项式系数为实数，非实根若出现必须带有共轭根；
唯一性因而还说明这个根是实数。
本节引用 Rouché 定理，但圆周界及其适用条件已逐项检查。

\section{找到一些零点与找全零点不同}
对一般复函数，要证明区域 $D$ 中的零点已找全，可以先由辐角原理计算总数，
再与各个隔离区域中的计数之和比较。若边界有零点，必须改变边界或单独处理，
不能直接使用无边界零点的公式。
对 $\zeta$ 而言，临界线上找到很多零点与临界带中已无遗漏，是不同任务；
Turing 型计数核对用于后一个任务，本篇不重新证明该方法；其逻辑结构见学习笔记第 16 篇。

\section*{参考依据与后续核验}
连续性、中值定理、二分法及 Newton 法的本例推导在文中直接给出；
Rouché 定理和辐角原理作为标准结果引用，可参看本书复分析部分。
关于 $\zeta$ 的认证背景可参看本书《非平凡零点的数值计算与显式公式逼近》。
本文尚未附具体区间计算软件的实现，避免把普通高精度浮点例程误标为认证程序；
区间运算的包含原理及用 Arb 球算术认证的实例，见学习笔记第 16 篇《计算机能证明定理吗》。
\end{document}
